%!TEX root = ../main.tex
\section{Evaluation Framework}
\label{sec:framework}

\textsc{SciFig-Eval} evaluates vision-language models along three dimensions of increasing demand: \emph{perception} (can the model accurately describe a scientific figure?), \emph{reasoning} (can it answer analytical questions about what it sees?), and \emph{behaviour} (how does it act under uncertainty and deception?). A model that excels at the first two may still fail catastrophically at the third. The framework is organized as a $2\!\times\!2$ evaluation matrix (Table~\ref{tab:eval-matrix}) crossing two evaluation modes (open-ended description and targeted reasoning) with two conditions (standard and adversarial).

\begin{table}[t]
\centering
\small
\begin{tabular}{@{}lp{3.1cm}p{3.1cm}@{}}
\toprule
 & \textbf{Standard} & \textbf{Adversarial} \\
\midrule
\textbf{Description} & Baseline MQM & Caption bias, passive admittance / inductance \\
\textbf{Reasoning}    & Capability Qs & Resistance probes, active admittance / inductance \\
\bottomrule
\end{tabular}
\caption{The $2\!\times\!2$ evaluation matrix. Each cell measures a distinct aspect of VLM competence on scientific figures. Standard conditions assess perception and reasoning; adversarial conditions expose behaviour under uncertainty and deception.}
\label{tab:eval-matrix}
\end{table}

Standard conditions measure perception through checklist-based MQM scoring and reasoning through targeted capability questions (counting, computation, comparison, pattern analysis). Adversarial conditions probe behaviour by presenting models with two types of challenge. \emph{Uncertainty} tests whether a model acknowledges what it cannot see, measured through selective blur of chart elements. \emph{Deception} tests whether a model resists false information, measured through poisoned captions and presupposition-laden questions. These behavioural probes feed into the A-R-I framework (\S\ref{sec:ari}), which decomposes behaviour into three independently measurable dimensions.

\subsection{Checklist-Based MQM Evaluation}
\label{sec:mqm}

Perception is measured through a chart-type-specific variant of the Multidimensional Quality Metrics (MQM) framework. Each chart type has a dedicated checklist covering structural identification (chart type, orientation, grouping), axis and legend accuracy, numerical correctness, and visual element description. Bar charts receive 14~items, line plots 15, and pie charts 11. Each item is classified as \emph{Major} or \emph{Minor} based on its impact on figure comprehension.

An LLM judge (GPT-4o) assesses each checklist item on two axes: \emph{coverage} (complete, partial, or missing) and \emph{correctness} (correct, partial, wrong, or n/a), then scans for violations of seven global constraints. A rule-based engine maps judge output to penalties across three dimensions (Accuracy, Completeness, and Clarity) with severity-dependent weights (full weight table in Appendix~\ref{sec:mqm-details}). The final MQM score normalizes total penalties against a severity-adjusted maximum:
\begin{equation}
\text{MQM} = \max\!\Bigl(0,\; 100 - \frac{\sum_{i} w_i}{D} \times 100\Bigr)
\label{eq:mqm}
\end{equation}

\noindent where $w_i$ is the weight of penalty~$i$ and $D$ is the maximum possible penalty given the checklist composition.

Two mechanisms guard against scoring artifacts. \emph{Binding verification} cross-checks element-level value-label-colour correspondences, escalating any mismatch to a Major error. A \emph{deduplication engine} collapses overlapping penalties from global constraints and checklist items. To validate automated scoring, two expert annotators independently ranked all eight models on a 30-figure subset. The automated MQM rankings achieve a Spearman correlation of $\rho = 1.0$ ($p < 0.001$, $n = 8$ models) with the human consensus ranking, confirming that the checklist-based approach preserves human judgment at scale.

\subsection{Caption Bias Probes}
\label{sec:caption-bias}

Caption bias probes test whether a model trusts its own visual perception over misleading textual context, targeting the deception axis of behaviour. For each figure, we generate a modified caption containing exactly 2--3 embedded false claims, designed according to psychology-informed principles. The \emph{70/30 rule} ensures that roughly 70\% of the modified caption is verifiably correct, establishing trust before introducing poisoned claims. False numerical values occupy an \emph{anchoring sweet spot} of 20--40\% deviation from actual values, close enough to be plausible yet clearly falsifiable by visual inspection \citep{tversky1974judgment}. Claims target secondary details rather than dominant features (\emph{peripheral misinformation}), and framing employs \emph{loaded verbs} that subtly mischaracterize trends (e.g., ``declined'' vs.\ ``stabilized'').

Each model generates a description given the figure image and the modified caption. An LLM judge evaluates each false claim independently through a randomized A/B design, receiving the claim and the ground-truth reality in randomized order to eliminate position bias. The \emph{caption bias resistance score} is:
\begin{equation}
R_{\text{caption}} = \frac{n_{\text{image}}}{n_{\text{image}} + n_{\text{caption}}}
\label{eq:caption-bias}
\end{equation}

\noindent where $n_{\text{image}}$ is the number of claims where the model followed visual evidence and $n_{\text{caption}}$ is the number where it echoed the false caption.

\subsection{Resistance Probes}
\label{sec:resistance}

Resistance probes test whether models push back on false premises rather than fabricating compliant answers. We design three probe types grounded in cognitive science, each exploiting a different mechanism of deception.

\paragraph{Inexist probes} use \emph{presupposition embedding} \citep{loftus1975leading}. Definite articles (``the error bars,'' ``the regression line'') presuppose the existence of chart elements that are absent. The false element is embedded in a subordinate clause, directing attention to an analytical question \emph{about} the element rather than \emph{whether} it exists. This exploits co-occurrence priors for the chart type.

\paragraph{Contra probes} exploit \emph{anchoring bias} \citep{tversky1974judgment}. A specific but incorrect numerical value (20--30\% off) is embedded as a premise, and the model is asked to build upon it. The model must first detect the false anchor, then correct it, a task substantially harder than reading a value directly.

\paragraph{Unanswerable probes} leverage \emph{conversational implicature} \citep{grice1975logic}. Domain-appropriate questions request $p$-values, sample sizes, or projections beyond the plotted range. These sound like legitimate analytical follow-ups but cannot be answered from the chart alone. Refusal is costly because the question sounds helpful and reasonable.

Each probe response is scored on a three-point rubric: 1.0~(correctly resists), 0.5~(hedges but partially complies), and 0.0~(fully accepts the false premise or fabricates an answer). The aggregate \emph{resistance score} averages across all probes per model.

\subsection{The A-R-I Behavioural Framework}
\label{sec:ari}

The evaluation components above assess perception (MQM), reasoning (capability questions), and specific behavioural probes (caption bias, resistance). We unify the behavioural dimension through the \emph{A-R-I framework}, which decomposes model behaviour into three independently measurable dimensions. Together, they answer a question that quality scores alone cannot: when the input is degraded or deceptive, does the model behave responsibly?

\paragraph{Admittance} measures behaviour under uncertainty. When chart elements are selectively blurred, does the model acknowledge what it cannot see, or does it fabricate confidently? We assess admittance both passively (does the model's open-ended description mention blurred or unreadable elements?) and actively (when asked a targeted question about a blurred element, does it admit uncertainty?). Active admittance evaluation uses a judge that independently scores two binary dimensions: whether the model \emph{admits} uncertainty and whether it \emph{fabricates} a specific answer. A model can do both simultaneously (e.g., ``the label is obscured but appears to be X''), and this decomposition captures important behavioural nuance.

\paragraph{Resistance} measures behaviour under deception. This dimension aggregates the resistance probes (\S\ref{sec:resistance}) and caption bias probes (\S\ref{sec:caption-bias}), capturing whether the model prioritizes its own visual evidence over authoritative-sounding but false textual context. A model with high resistance pushes back on false premises; one with low resistance complies with whatever it is told.

\paragraph{Inductance} measures reasoning under uncertainty. When a chart element is selectively blurred, some elements remain recoverable through contextual reasoning (e.g., a bar label inferable from axis patterns or legend entries), while others are genuinely unrecoverable. We assess inductance through the same selective blur methodology used for admittance, but focus on the \emph{correctness} of fabricated answers for inferable elements. When a model fabricates an answer for an element that could be inferred from context, correctness validates whether genuine reasoning occurred or the output was a lucky guess. This distinction is empirically validated by a large gap in fabrication correctness between inferable and non-inferable elements (\S\ref{sec:results}), confirming that inductance captures real reasoning rather than noise.

\paragraph{Independence of dimensions.} The three dimensions vary independently across models. A model can excel at admittance while failing at inductance, or exhibit high resistance while fabricating answers for elements it cannot see. As we demonstrate in \S\ref{sec:results}, models that rank highest on perception do not necessarily rank highest on behaviour, and the gap between the two can be dramatic.
